\scanpage{41}
% L1–5
\deleted{\deleted{\unsure{sur|sont}} en un pt \unsure{au} dessus de $s$, i.e. de $X_0$. Donc pour voir que c'est un isom \unsure{en} pt \unsure{gén.|géom.} $\unsure{z|\bar z|Z}$ de $X_0$, il suffit de voir que \deleted{\unsure{c'est}} \inserted{en} lui appliquant $\unsure{\mathbb R|R}j^!$, où $j:X_0\to Y$, on trouve un isomorphisme en $z$. Or on a}

% L6–7
\diagram{\begin{tikzcd}[ampersand replacement=\&,column sep=large,row sep=small]
\deleted{\mathbb Rj^!(\mathbb Ri^!(\mu_{\alpha}))}\arrow[r]\arrow[d,phantom,"\deleted{\unsure{\wr\!\vert|\Vert|\wr}}" description]\&\deleted{\mathbb Rj^!(\unsure{\Lambda|A}_Y[-2]]}\arrow[d,phantom,"\deleted{\Vert}" description]\\
\deleted{\mathbb Rh^!\deleted{\unsure{i^!}}(\unsure{\mu_X|N_X|\mu_\alpha})}\&{}
\end{tikzcd}}

% L7, left
\diagram{\begin{tikzcd}[ampersand replacement=\&,column sep=small]
\deleted{X_0}\arrow[r,"\deleted{j}"]\arrow[rr,bend right=23,dash,"\deleted{h}"']\&\deleted{Y}\arrow[r,"\deleted{i}"]\&\deleted{X}
\end{tikzcd}}

% archive sheet number: 40
